\begin{problem}[The adjunction in full natural form]
Let $a:\PreSh_X\to\Sh_X$ be sheafification and $i:\Sh_X\to\PreSh_X$ inclusion. Construct explicitly the bijection
\[
\Phi_{\mathcal F,\mathcal G}:\Hom_{\Sh_X}(a\mathcal F,\mathcal G)\to\Hom_{\PreSh_X}(\mathcal F,i\mathcal G)
\]
and prove it is natural in both variables.
\end{problem}

\paragraph{Why this problem matters.}
Adjunction packages the universal property into the categorical statement used throughout later homological sheaf theory.

\paragraph{Useful ideas before starting.}
Use the universal property, stalk preservation, and the fact that a morphism of sheaves is an isomorphism when it is so on every stalk.  Keep the distinction between a presheaf statement on sections and a local statement on germs explicit.

\paragraph{Guided hints.}
The bijection is precomposition with the unit. Naturality follows by associativity plus the defining naturality of the unit.

\ifdefined\FullProblemDossiers
\begin{solution}
Set
\[
\Phi_{\mathcal F,\mathcal G}(h)=i(h)\circ\eta_{\mathcal F}.
\]
The universal property says exactly that this map is bijective. For $u:\mathcal F'\to\mathcal F$ and $v:\mathcal G\to\mathcal G'$ with sheaf targets, associativity gives
\[
\Phi_{\mathcal F',\mathcal G}(h\circ a(u))
=h\,a(u)\eta_{\mathcal F'}=h\eta_{\mathcal F}u
=\Phi_{\mathcal F,\mathcal G}(h)\circ u,
\]
using naturality $a(u)\eta_{\mathcal F'}=\eta_{\mathcal F}u$. Similarly
\[
\Phi_{\mathcal F,\mathcal G'}(v\circ h)=v\circ\Phi_{\mathcal F,\mathcal G}(h).
\]
Thus the bijections are natural in both variables and $a\dashv i$.
\end{solution}

\paragraph{What happened mathematically.}
The adjunction is the universal property with both variables allowed to vary.

\paragraph{Alternative approaches.}
Whenever possible one may replace the espace-etale model by the two-plus construction, or conversely replace a cover/refinement argument by the universal property.  The two approaches agree because both construct the same associated sheaf.

\paragraph{Extensions.}
Use left-adjointness to deduce preservation of colimits.

\paragraph{Related problems.}
Compare with VI.14.P1 for the complete legacy construction and with the later sheaf chapters for operations that are deliberately not migrated here.

\paragraph{Provenance.}
Legacy INCLUDE-derived: the adjunction demand is an explicit final clause of AG1 Exercise 4 and remains fully solved in P1.
\else
\paragraph{Provenance.}
Legacy INCLUDE-derived: the adjunction demand is an explicit final clause of AG1 Exercise 4 and remains fully solved in P1.
\fi
